\documentclass{situationist-preprint}

\title{Expertise and Authority}
\msid{POL-15-b}
\msfamily{POL-15}
\msdiscipline{Political Philosophy}
\mstype{Consequence}
\mskeywords{expertise; authority; epistocracy; delegation; plural voting; technocracy}

\begin{document}
\maketitle

\begin{abstract}
POL-15-a argues that philosophical expertise does not confer political authority: knowledge is not authority, and the one route from the first to the second, reasonable deference, is closed where the identification of experts is subject to reasonable disagreement and deference cannot be confined to the expert's domain of reliability. We derive two consequences. First, the argument generalizes from philosophers to any claimed expertise on questions of justice, and so excludes epistocratic schemes such as Mill's plural voting. Second, it does not exclude the authority of technical bodies such as central banks or drug regulators, because there the two conditions that close the route can fail; but the authority of such bodies is delegated, and derives from a democratic source.
\end{abstract}

\section{Two consequences}

The argument of POL-15-a has two steps. Proposition~2.3 states the gap: from the fact that $E$ knows better than $S$ what ought to be done it does not follow that $E$ has authority over $S$. Proposition~2.4 states that the route across the gap, that it would be reasonable for $S$ to defer to $E$, is closed in the political case, for two reasons: (a) who the experts on justice are is itself subject to reasonable disagreement, and (b) deference on questions of justice cannot be confined to the questions on which the putative expert is reliable. Neither step mentions philosophers except as the subject of the claim.

\begin{proposition}\label{p:epistocracy}
No claimed expertise on questions of justice, philosophical or otherwise, confers political authority. In particular, schemes that weight citizens' votes by education or competence lack authority over those whose votes are discounted.
\end{proposition}

\begin{proof}
Conditions (a) and (b) concern the questions, not the class of expert. Any scheme that assigns greater political power on the basis of competence must identify the competent, and the criteria of identification for questions of justice are subject to reasonable disagreement \citep{Estlund2008}; and the discounted citizen cannot confine her deference to the questions on which the favored are reliable, since the scheme gives them weight on every question decided by vote. By Proposition~2.4 and Theorem~2.5 of POL-15-a, the scheme has no authority over her.
\end{proof}

The proposition applies to Mill's proposal that the educated have additional votes \citep[ch.~8]{Mill1861} and to its recent epistocratic successors \citep{Brennan2016}. It does not show that such schemes would produce worse decisions, which is an empirical question; it shows that they would lack the authority their decisions would claim.

\begin{proposition}\label{p:delegation}
Technical bodies can have authority within a mandate, but only by delegation from a body whose authority does not rest on expertise.
\end{proposition}

\begin{proof}
For a central bank setting interest rates within an inflation target, or a regulator approving medicines against a standard of safety, conditions (a) and (b) can fail: the relevant expertise is identified by criteria not reasonably disputed, and the mandate confines the body to the questions on which it is reliable. The route of reasonable deference is then open, as in Raz's normal justification thesis \citep{Raz1986}. But the mandate itself, the target and the standard, settles questions of justice and priority on which (a) and (b) hold. Its authority must therefore come from elsewhere; on the view of POL-04, from democratic procedures.
\end{proof}

\begin{objection}
If technical expertise can ground authority within a mandate, so can philosophical expertise within a suitably narrow one.
\end{objection}

\begin{reply}
It can, if a mandate could be drawn that confined the philosopher to questions on which her reliability is not reasonably disputed: on whether an argument is valid, for instance. Such a mandate would give her no authority over what is just, which was the claim of the philosopher-king.
\end{reply}

\section{Conclusion}

The argument against philosopher-kings is an argument against rule by experts on justice of any kind, and it is compatible with expert bodies exercising delegated authority on technical questions. Experts may advise, and may administer what a democratic authority has decided; this is the consequence of POL-15 for the place of expertise in government.

\begin{thebibliography}{9}\small
\bibitem[Brennan(2016)]{Brennan2016} Brennan, J. (2016). \emph{Against Democracy}. Princeton University Press.
\bibitem[Estlund(2008)]{Estlund2008} Estlund, D. M. (2008). \emph{Democratic Authority: A Philosophical Framework}. Princeton University Press.
\bibitem[Mill(1861)]{Mill1861} Mill, J. S. (1861). \emph{Considerations on Representative Government}. Parker, Son, and Bourn.
\bibitem[Raz(1986)]{Raz1986} Raz, J. (1986). \emph{The Morality of Freedom}. Clarendon Press.
\end{thebibliography}

\end{document}
